\chapter{Discussion}
\label{ch:discussion}

This chapter (1) interprets the experimental results of
Chapter~\ref{ch:experiments}, then distills the engineering experience gained
while building and operating the pipeline into lessons on (2) exploring live
websites, (3) guarding LLM agents structurally rather than through prompt
wording, (4) inducing high-quality planning operators, and (5) training and
serving small language models, before (6) summarizing these lessons as best
practices and (7) discussing threats to validity. Each lesson follows the
same pattern: a problem actually encountered in our runs, its diagnosis, the
fix, and the general principle we draw from it. Full incident details are
collected in Appendix~\ref{app:runs}.

\section{Interpreting the Results}
\label{sec:disc-results}

% TODO: refresh this section once the remaining three sites and the
% statistical repeats (Sections 4.5-4.6) are in.

The headline number so far --- an improvement from 5.56\% to 6.67\% task
success on the official GitLab suite --- deserves a careful reading rather
than a triumphant one. In absolute terms the gain is small, and with one
attempt per task it must be treated as a directional signal until the
repeated runs of Section~\ref{sec:exp-stats} are complete. What makes the
signal interesting is its cost profile: it was obtained from only 352
training examples, collected autonomously from a website with no human
annotation, and the fine-tuning itself took under half an hour on a single
GPU. The pipeline, in other words, moved a small model in the right direction
with a remarkably small amount of data; the open question is how far the
trend continues as data collection scales, which the cross-site results of
Section~\ref{sec:exp-main} will begin to answer.

Three factors plausibly bound the current gain. First, dataset size: 352
examples is two to three orders of magnitude below typical agent fine-tuning
corpora, so the model sees each interaction pattern only a handful of times.
Second, coverage: the training data was collected on a subset of site areas,
while the benchmark draws tasks from the whole site. Third, the base model
itself: a 9B-parameter model fine-tuned with supervised trajectories alone
cannot be expected to recover capabilities --- long-horizon planning,
self-correction --- that the trajectories only sparsely demonstrate. None of
these factors indicts the pipeline design; all three point to data volume
and breadth as the binding constraint, which is precisely the constraint the
pipeline is built to relax.

Equally important is what the end-to-end run \emph{demonstrated}: every stage
--- exploration, discovery, operator chaining, curation, training, serving,
and official evaluation --- executed against a real website without manual
intervention, and the quality gates rejected data that would otherwise have
polluted training (Section~\ref{sec:disc-slm}). The remainder of this chapter
reports what it took to make that sentence true.

\section{Lessons from Exploring Live Websites}
\label{sec:disc-exploration}

\subsection{Frontier starvation}

\textbf{Problem.} An early exploration run on the Reddit-style site stalled
at 57 of a budgeted 500 nodes: the frontier simply ran dry.
\textbf{Diagnosis.} The frontier then grew only from URLs the solver happened
to navigate to. Pages whose proposed tasks were all informational produced no
navigation at all, so entire regions of the site were never queued ---
exploration was hostage to which tasks the explorer happened to propose.
\textbf{Fix.} Hyperlink mining was added as a second frontier tier, with URL
templating and per-template caps to keep near-duplicate pages from flooding
the queue (Section~\ref{sec:method-frontier}).
\textbf{Lesson.} On a real website, \emph{reachability must not depend on
task-driven behavior}. Structural signals (links) and behavioral signals
(navigation) are complementary frontier sources, and the structural one needs
explicit diversity controls because it arrives in bulk.

\subsection{Task proposal collapse}

\textbf{Problem.} In one production run, 46\% of all proposed tasks were
near-identical across pages --- variations of ``go to the Forums page''
re-proposed on nearly every node.
\textbf{Diagnosis.} Every page of a site shares its global navigation menu,
and an explorer prompted page-by-page has no memory of what it already
proposed elsewhere.
\textbf{Fix.} The explorer now persists for the whole run: recent proposals
are echoed into its prompt as negative context, and near-duplicates are
additionally dropped in code (Section~\ref{sec:method-lowlevel}).
\textbf{Lesson.} Diversity of generated tasks is a \emph{run-level} property,
not a page-level one; it must be enforced by state that outlives any single
LLM call.

\subsection{Unreliable pages}

\textbf{Problem.} A single slow page --- a navigation timeout on one mined
URL --- crashed an entire multi-hour exploration run.
\textbf{Diagnosis.} Browser operations were implicitly trusted; one
unhandled exception anywhere in node processing propagated to the top level.
\textbf{Fix.} Every browser operation is now fail-soft: a node that cannot be
opened is logged, recorded with its error, and skipped, while genuinely
unexpected exceptions still abort the run loudly. Artifacts are persisted
after every node, so even an aborted run leaves a consistent, resumable
snapshot.
\textbf{Lesson.} Treat every interaction with a live website as an
operation that \emph{will} eventually fail, and decide per operation whether
failure is survivable; make the run's progress durable at the finest
granularity the budget allows.

\subsection{Shared, mutable state}

\textbf{Problem.} Benchmark websites are shared and stateful: a careless
agent can pollute --- or destroy --- the environment other runs depend on.
\textbf{Fix.} Three layers of discipline: a read-only mode that drops
state-modifying tasks fail-closed (any task the explorer fails to classify is
treated as modifying); a standing rule forbidding account-level destructive
tasks regardless of mode; and, for evaluation on a live shared instance, an
intent-based classification of benchmark tasks into read-only and
write-mutating subsets so that two machines can sweep the same site
concurrently without interfering (Section~\ref{sec:disc-slm}).
\textbf{Lesson.} On shared infrastructure, safety must be
\emph{fail-closed}: the burden of proof lies on a task being harmless, not on
it being harmful.

\section{Structural Guardrails over Prompt Nudges}
\label{sec:disc-guardrails}

A recurring theme of our runs is that instructing a model to behave is not
the same as making the system behave. The incidents below each began as a
prompt-level convention that a real model quietly violated, and each ended as
a structural mechanism that does not depend on model obedience.

\subsection{Tag conventions do not survive contact with real models}

\textbf{Problem.} Task proposals were originally classified by a leading
\texttt{[INFO]}/\texttt{[NAV]}/\texttt{[MOD]} tag that the explorer was
instructed to emit. In a real run, 89\% of proposals carried no tag at all.
\textbf{Diagnosis.} The model's natural tendency was to return the
classification as a structured \texttt{type} field --- which the parser
discarded --- rather than as a bracket inside a string. Worse, the safety
check for modifying tasks was a literal substring test on \texttt{[MOD]}, so
an untagged destructive task would have sailed through it.
\textbf{Fix.} Proposals are now requested and parsed as structured objects
with an explicit type field; anything unclassifiable is tagged
\texttt{UNKNOWN} and dropped in read-only mode.
\textbf{Lesson.} Ask models for \emph{structure}, not for formatting
conventions inside free text --- and never build a safety check on a
convention the model may not follow.

\subsection{The stuck loop, in two iterations}

\textbf{Problem.} Solvers wasted entire step budgets repeating an action that
visibly did nothing --- in one recorded episode, clicking the same inert
element for all fifteen steps.
\textbf{Diagnosis and first fix.} We added a warning, injected into the
prompt after two consecutive no-effect actions on the same target. The next
real run falsified the fix: the model kept clicking with the warning present
from the third call onward. It also revealed a subtlety --- the model
sometimes retried a stuck dropdown with different-looking but equally wrong
arguments, so matching identical action strings would miss the loop entirely.
\textbf{Second fix.} The loop detector now keys on the \emph{target} acted
on and on whether the observation changed; after a bounded number of
no-progress repeats, the episode is terminated outright, spending no further
LLM calls. Requiring ``no observable change'' keeps the guard from misfiring
on legitimate sequences that reuse one element across changing pages.
\textbf{Lesson.} A prompt warning is a suggestion; a budget is a mechanism.
When a model demonstrably ignores the suggestion, escalate to a mechanism ---
and validate the escalation on real behavior, not on the behavior the prompt
requests.

\subsection{Smaller incidents, same shape}

Three further incidents followed the same pattern and received the same style
of fix. Explorer-proposed information tasks sometimes contained their own
answers (``find the post titled X''), which a solver can ``solve'' without
learning anything; generation now instructs against it and leaky tasks are
filtered. Vision-attached calls occasionally returned empty responses
(upstream refusals); the agent retries once in text-only mode rather than
failing the step. And episodes that exhausted their budget originally ended
in an uninformative failure; they now receive one final text-only wrap-up
call whose \emph{only} admissible outputs are a completion claim or an
infeasibility report --- turning a silent timeout into a labeled outcome that
the dataset can use.

\section{Lessons from Operator Induction}
\label{sec:disc-operators}

The semantic-task mechanism of Section~\ref{sec:method-semantic} is only as
good as the operators it chains. Two failure modes surfaced when we audited
the operators induced in an early production run, and both led to structural
gates rather than prompt adjustments.

\subsection{Tautological operators}

\textbf{Problem.} Three of the four operators induced in that run had effect
sets identical to their precondition sets.
\textbf{Diagnosis.} Asked to describe a skill's preconditions and effects,
an LLM finds it easiest to restate the context --- ``user is on the issue
list'' before, ``user is on the issue list'' after. Such an operator encodes
no state change, so it can never appear mid-chain and never lead anywhere
new; left unchecked, tautologies would have filled the operator graph with
dead nodes.
\textbf{Fix.} The anti-tautology gate: effects that restate a precondition
are dropped, and a skill left without novel effects is rejected from
promotion altogether, with the rejection logged.
\textbf{Lesson.} For LLM-authored symbolic artifacts, \emph{define validity
set-theoretically and enforce it in code}; do not rely on the prompt to
forbid degenerate answers.

\subsection{Synonymous predicates that embeddings cannot unify}

\textbf{Problem.} Chaining requires an effect of one operator to match a
precondition of another, but independently induced operators phrase the same
condition differently.
\textbf{Diagnosis.} We measured the embedding similarity of predicate pairs
from the same run: a true paraphrase pair scored 0.570 --- \emph{below} the
0.6 canonicalization threshold --- while a semantically different pair scored
0.606, above it. No threshold separates these cases; purely geometric
canonicalization cannot recover a shared vocabulary after the fact.
\textbf{Fix.} Prevention rather than cure: induction is vocabulary-aware.
The prompt presents the registry's existing predicate labels with an
instruction to reuse an equivalent one verbatim, so unification happens at
authoring time, where the model has full context, instead of at embedding
time, where only the surface strings remain.
\textbf{Lesson.} When two LLM outputs must later \emph{match}, share the
vocabulary at generation time; do not defer semantic unification to a
similarity metric.

\section{Pitfalls in Training and Serving SLMs}
\label{sec:disc-slm}

\subsection{The chat template is part of the model}

The target model family emits explicit reasoning segments by default, while
our training targets are bare thought--action pairs. The correct,
thinking-disabled chat template had to be identified from the training
framework's template registry rather than guessed from documentation, and the
same setting had to be applied again, per request, at serving time. A
mismatch on either side silently trains or evaluates a different
distribution than intended. The general point: \emph{template and decoding
configuration are part of the model artifact}, and train-time and
serve-time configurations must be pinned together.

\subsection{The prompt-fidelity bug}

The most instructive failure of the training phase produced no error at all.
A quick post-training smoke test, written with an \emph{approximation} of the
solver prompt, made the fine-tuned model emit natural-language descriptions
instead of executable actions --- which looked exactly like a failed
fine-tune. The model was fine; the test prompt was unfaithful. Re-testing
with the production prompt --- imported from the solver, not re-typed ---
produced correct actions. This incident is why prompt parity
(Section~\ref{sec:method-data}) is enforced by \emph{importing} the live
system prompt and its helper routines into the dataset builder and the
evaluation harness, never by duplicating them. Any test of a fine-tuned
agent is only as valid as the fidelity of its prompt to production.

\subsection{Dependency isolation}

Supporting a newly released model architecture required the development
version of the modeling library, which conflicted with the serving engine's
pinned requirements. Three isolated environments --- training, serving, and
evaluation --- resolved what no single environment could satisfy. Mundane,
but it converts an unsolvable version conflict into a routine setup step.

\subsection{Agents grade their own homework}

The judge pass exists because declared success is not success: a solver that
answers ``done'' has evaluated itself. Re-checking declared successes against
the episode's own final state caught self-graded failures before they became
training examples --- and the dataset builder's strict mode inverts the
default, dropping any declared success that was never judged rather than
trusting it. Under principle P1, \emph{unverified data is not neutral; it is
a liability that compounds through training}.

\subsection{Evaluating safely on shared live sites}

Running baseline and fine-tuned sweeps against one live site instance would
normally serialize a full working day of evaluation. Classifying benchmark
tasks by intent into read-only and write-mutating subsets let a second
machine evaluate the read-only subset concurrently with the first machine's
unrestricted sweep. The classification itself held two traps worth recording:
noun-phrase collisions (``merge request'' contains a verb only
syntactically; ``report bugs'' does not name the reporting action the regex
expected) and word insertion defeating naive phrase patterns (``open a
\emph{new} issue''). Transfer cost also mattered: moving the 1.7\,GB LoRA
adapter and re-merging locally beat moving an 18\,GB merged checkpoint.

\section{Best Practices}
\label{sec:disc-practices}

Table~\ref{tab:practices} condenses the lessons of this chapter into
practices we would adopt from the first day of a similar project.

\begin{table}[H]
\centering
\caption[Distilled best practices]{Distilled best practices for
LLM-agent data-collection pipelines on live websites.}
\label{tab:practices}
\small
\begin{tabular}{L{4.4cm} L{8.8cm}}
\toprule
\textbf{Practice} & \textbf{Rationale (section)} \\
\midrule
Verify, then admit & Skills, operators, and successes are liabilities until independently checked (\ref{sec:disc-slm}, \ref{sec:disc-operators}). \\
Prompt parity by import & Training data, live agent, and eval harness share one prompt source; duplication drifts (\ref{sec:disc-slm}). \\
Structure over conventions & Request typed fields from models; never parse safety-relevant information out of free text (\ref{sec:disc-guardrails}). \\
Mechanisms over warnings & Convert repeated model misbehavior into budgets and hard stops, validated on real runs (\ref{sec:disc-guardrails}). \\
Fail-soft browsing & Any page operation may fail; skip and log locally, crash only on the truly unexpected (\ref{sec:disc-exploration}). \\
Persist per node & Interrupted runs should leave consistent, resumable artifacts (\ref{sec:disc-exploration}). \\
Fail-closed safety & Unclassifiable tasks count as destructive; destructive account actions are banned outright (\ref{sec:disc-exploration}). \\
Share vocabulary at authoring time & Symbolic artifacts that must match later should be generated against a shared label set (\ref{sec:disc-operators}). \\
Pin template with model & Chat template and decoding flags are part of the artifact, for training and serving alike (\ref{sec:disc-slm}). \\
Log everything as events & A per-run event log makes post-hoc failure analysis --- most of this chapter --- possible at all (\ref{sec:disc-exploration}). \\
\bottomrule
\end{tabular}
\end{table}

\section{Threats to Validity}
\label{sec:disc-validity}

\textbf{Internal validity.} The evaluation to date runs each benchmark task
once per model; individual task outcomes on a live site are noisy, so the
measured gap could shift under repetition. The planned repeated runs
(Section~\ref{sec:exp-stats}) address this directly. Several quality gates
(the judge, the re-annotator) are themselves LLMs and inherit LLM biases;
we mitigate but do not eliminate this by checking judgments against concrete
final states rather than against transcripts alone.

\textbf{External validity.} All data and evaluation come from self-hosted
WebArena sites --- realistic, but four sites of one benchmark family ---
and from one 9B model. Claims are correspondingly scoped: we report that the
pipeline improves \emph{this} model on \emph{these} sites, and treat broader
generalization as future work rather than as an implication.

\textbf{Construct validity.} Task success against official evaluators is the
field-standard construct, but it is binary and end-state-based: it does not
credit partial progress, and semantic-task episodes verified by self-report
plus offline judging are a weaker construct than programmatically checked
ones. We keep the two verification regimes distinguishable in the data so
that stricter analyses remain possible.

\vspace{1em}
\noindent
These lessons, and the results they contextualize, complete the empirical
story of the thesis. The final chapter summarizes what was built and learned,
and where the work goes next.
